\section{Kalman 滤波：用预测与更新的循环从噪声中估计状态}
\label{sec:16-Control-and-Reinforcement-Learning/01-Classical-Control/04}

你控制着一间冷库。温度传感器的读数在变——一部分是因为压缩机真的在制冷，一部分是因为有人开门卸货，还有一部分纯粹是传感器的电噪声。你的控制器需要当前的温度状态来做决策，但你拿到的只有被噪声污染的读数。真实的温度躲在传感器后面，你看不到它。

这不是一个数据平滑问题。你不能等到收齐了今天全部读数再回头算——控制器需要在每个时刻立刻给出当前的最佳估计，然后马上拿这个估计去算下一拍的控制指令。你要的是一个在线递推：新数据来了，更新一次估计，扔掉旧数据，不必重新遍历历史。

这就是状态估计问题。Kalman filter 用一套简洁的预测-更新循环解决了它，前提是系统是线性的、噪声是 Gaussian 的。但这套前提一旦不满足，filter 给出的估计可能不只是不准，而是系统性地偏到错误的地方去。我们先从它怎么工作讲起，最后再划出它不再可靠的那条线。

\subsection{一个具体的困境：传感器读数不等于状态}

冷库的简化热平衡方程离散化后可以写成：

\[
T_k = a T_{k-1} + b u_k + w_k,
\]

其中 \(T_k\) 是时刻 \(k\) 的仓内温度，\(u_k\) 是压缩机控制量，\(w_k\) 吸收了墙体蓄热、开门和货物进出中无法用方程精确描述的部分——把它建模为零均值 Gaussian 噪声。

你拿到的传感器读数 \(z_k\) 不等于 \(T_k\)。传感器可能有一个缓慢漂移的偏置 \(b_k\)，外加短时的随机电噪声 \(\varepsilon_k\)：

\[
z_k = T_k + b_k + \varepsilon_k.
\]

如果模型把偏置 \(b_k\) 固定为零，filter 会越来越相信自己的温度估计，而对传感器反馈越来越不敏感——最终那个偏置把你带离真实温度，filter 却自信满满地报告一个错误的值。

把偏置也放进状态向量，令 \(x_k = [T_k, b_k]^{\trans}\)，并假设偏置按随机游走演化 \(b_k = b_{k-1} + \zeta_k\)，Kalman filter 就能同时估计温度和偏置。但温度与偏置能否从现有的输入和观测中可靠地分开，取决于模型的结构性质——后文的可观测性——而不是递推步数够不够多。

\subsection{两个高斯，一次联合：更新步骤的心跳}

先写出一般形式的线性 Gaussian 状态空间模型：

\[
X_t = A X_{t-1} + B u_t + W_t,\qquad W_t \sim \mathcal{N}(0, Q),
\]
\[
Y_t = C X_t + V_t,\qquad V_t \sim \mathcal{N}(0, R),
\]

并假设过程噪声、观测噪声与初始状态之间适当独立。

在线性与 Gaussian 的双重假设下，一个关键的闭包性质成立：如果上一时刻的后验是 Gaussian，那么预测分布、联合分布以及更新后的后验全都是 Gaussian。整条递推只需要维护两个量——均值向量和协方差矩阵——不需要追踪整个分布的形状。

假设我们在 \(t-1\) 时刻已经拥有

\[
X_{t-1} \mid y_{1:t-1} \sim \mathcal{N}(m_{t-1},\, P_{t-1}).
\]

\textbf{预测步骤}：把上一时刻的后验推过动力学。线性变换加独立 Gaussian 噪声，均值和协方差的传播有闭式：

\[
m_t^- = A m_{t-1} + B u_t,
\]
\[
P_t^- = A P_{t-1} A^{\trans} + Q.
\]

上标 \(-\) 表示``在看见 \(y_t\) 之前''。即使均值按确定性动力学演化，过程噪声 \(Q\) 也会增加协方差——它在说：动力学本身有未建模的部分，时间往前推一步，不确定性的基线就要往上加一截。如果你在实现中忘了加 \(Q\)，filter 会越来越自信，最终拒绝真实观测。

\textbf{更新步骤}：现在 \(y_t\) 到达。你把预测状态和新观测放在一起看它们的联合分布：

\[
\begin{bmatrix}
X_t \\ Y_t
\end{bmatrix}
\sim
\mathcal{N}\!\left(
\begin{bmatrix}
m_t^- \\ C m_t^-
\end{bmatrix},
\begin{bmatrix}
P_t^- & P_t^- C^{\trans} \\
C P_t^- & C P_t^- C^{\trans} + R
\end{bmatrix}
\right).
\]

这个联合协方差矩阵的下三角块 \(C P_t^-\) 是预测状态与预测观测之间的协方差——它度量了``如果真实状态比预测的高一点，观测值也会高多少''。

定义 innovation（新息）：

\[
e_t = y_t - C m_t^-,
\]

即实际观测与预测观测之间的偏差。定义 innovation 协方差：

\[
S_t = C P_t^- C^{\trans} + R.
\]

对联合 Gaussian 应用条件分布公式——给定 \(Y_t = y_t\) 时 \(X_t\) 的条件分布仍然是 Gaussian，其均值和协方差为：

\[
K_t = P_t^- C^{\trans} S_t^{-1},
\]
\[
m_t = m_t^- + K_t e_t,
\]
\[
P_t = P_t^- - K_t C P_t^-.
\]

\(K_t\) 叫 Kalman gain，它是这一整套推导的核心输出。把它拆开来看：\(P_t^- C^{\trans}\) 是状态与观测的协方差——``状态的哪个方向在观测中有投影''；\(S_t^{-1}\) 是观测不确定性的逆——``观测有多可信''。二者相乘的结果决定了：在观测告诉你的方向上，你用多大的步长修正预测。

\begin{itemize}
\tightlist
\item
  观测噪声 \(R\) 大（传感器很吵）\(\rightarrow\) \(S_t\) 大 \(\rightarrow\) \(K_t\) 小 \(\rightarrow\) 更信任动力学预测，不怎么信新读数。
\item
  预测不确定性 \(P_t^-\) 大，且观测能看到相应方向 \(\rightarrow\) \(K_t\) 大 \(\rightarrow\) 新数据大幅修正估计。
\item
  观测噪声 \(R\) 几乎为零，预测不确定性不为零 \(\rightarrow\) \(K_t \approx C^{-1}\)（若 \(C\) 可逆），估计几乎完全跟随观测。
\end{itemize}

Kalman gain 不是在两个固定权重之间做选择——它会随每一步的 \(P_t^-\) 动态变化。filter 刚启动时，初始协方差很大，gain 很大，快速吸收早期数据。稳定运行后，\(P_t\) 收敛到稳态值，gain 也趋于常数。

\subsection{走一个小规模数值例子}

设一维位置跟踪问题。状态只有位置 \(x_t\)，动力学是随机游走：\(x_t = x_{t-1} + w_t\)，\(w_t \sim \mathcal{N}(0, 0.1)\)。观测直接看位置加噪声：\(y_t = x_t + v_t\)，\(v_t \sim \mathcal{N}(0, 1)\)。

初始估计：\(m_0 = 0\)，\(P_0 = 10\)（对初始位置高度不确定）。

第一步预测：\(m_1^- = 0\)，\(P_1^- = 10 + 0.1 = 10.1\)。

假设施加控制输入后第一次真实观测 \(y_1 = 1.2\)。计算 Kalman gain：

\[
K_1 = \frac{10.1}{10.1 + 1} \approx 0.91.
\]

更新：

\[
m_1 = 0 + 0.91 \times (1.2 - 0) \approx 1.09,
\]
\[
P_1 = 10.1 - 0.91 \times 10.1 \approx 0.91.
\]

协方差从 10.1 骤降到 0.91——第一次观测带来了大量信息。gain 0.91 说明 filter 几乎完全信任这次观测（观测噪声 1 相对于预测不确定性 10.1 确实很小）。

再做一步。预测：\(m_2^- = 1.09\)，\(P_2^- = 0.91 + 0.1 = 1.01\)。假设 \(y_2 = 1.5\)。gain：\(K_2 = 1.01 / (1.01 + 1) \approx 0.50\)。更新：\(m_2 \approx 1.30\)，\(P_2 \approx 0.50\)。gain 降为 0.50——经过第一步，filter 已经比较确定了，对后续观测的采纳更谨慎。

几步之后，\(P_t\) 将趋于稳态值。求解代数 Riccati 方程 \(P = A P A^{\trans} + Q - A P C^{\trans}(C P C^{\trans} + R)^{-1} C P A^{\trans}\) 可以得到稳态协方差的闭式解。在这个例子里稳态 \(P\) 大约在 0.3 附近——无论跑多少步，协方差不会再低于它，因为每一步的过程噪声都在持续注入新的不确定性。

这个例子暴露了 Kalman filter 的权衡本质：过程噪声 Q 和观测噪声 R 的相对大小决定了 filter 最终信任动力学和信任观测之间的平衡点。如果 Q 设得太小，filter 对动力学过于自信，观测修正的效果被压制，估计会滞后于真实状态的变化——这称为估计滞后（estimator lag），是手动调低 Q 最常见的副作用。

\subsection{可观测性：有些方向再多数据也看不见}

Kalman filter 的协方差 \(P_t\) 会缩小，但不是所有方向都缩得一样快。某个状态方向如果长期不通过 \(C X_t\) 影响观测，filter 就收不到关于它的信息。动力学 \(A\) 可以把该方向的能量逐步耦合到可观测方向——但也可能永远不耦合，让那个方向一直藏在暗处。

线性系统的 observability 矩阵

\[
\mathcal{O} =
\begin{bmatrix}
C \\ CA \\ \vdots \\ CA^{d-1}
\end{bmatrix}
\]

回答了这个问题：如果 \(\mathcal{O}\) 满列秩，在理想无噪条件下，有限时间段内的观测足以区分任意两个不同的初始状态。如果不满秩，则存在不可观测的子空间——其中的状态分量无论采集多少数据都无法从观测中恢复，即使 \(R\) 任意小。

这一点对实践者不是纸上谈兵。回到冷库的例子：如果温度传感器只有一个，而状态包含了温度和偏置，observability 矩阵能否满秩取决于动力学 \(A\) 中温度与偏置是否有足够的耦合路径。如果偏置漂移极其缓慢且温度演化与其无关，两个状态的分离可能需要极长的时间——在有限时间窗内 filter 给出的偏置估计可能毫无意义。不是递推算错了，是信息根本就不在数据里。

对偶的概念是可控性——通过控制输入 \(u_t\) 能否把状态驱动到任意目标，取决于 controllability 矩阵的秩。可控性和可观测性一起构成了线性系统理论的骨架：可控性保证你能把系统带到想去的地方，可观测性保证你知道系统当前在什么地方。

\subsection{Filtering、prediction、smoothing 不是同一件事}

这三种任务经常被混用，但它们在使用数据的权限上有根本区别：

\begin{itemize}
\tightlist
\item
  Filtering：\(p(x_t \mid y_{1:t})\)——只使用当前及过去观测。这是在线控制的唯一合法选择。
\item
  Prediction：\(p(x_{t+h} \mid y_{1:t})\)——从当前往后推 \(h\) 步，不使用未来的任何信息。
\item
  Smoothing：\(p(x_t \mid y_{1:T})\)，其中 \(T > t\)——整段数据到齐后回头修正过去的状态估计。
\end{itemize}

离线轨迹重建可以放心用 smoothing——你已经看到全程数据了，允许未来观测修正你对过去的判断。但实时控制不能在时刻 \(t\) 使用尚未发生的 \(y_{t+1:T}\)。如果评价模型时误把平滑后的状态当成在线特征来用，就会出现时间泄漏：模型在训练时偷看了未来，上线后做不到同样的精度。

\subsection{数值实现不是套公式就行了}

写在纸上的公式用 \(S_t^{-1}\) 是为了表达线性算子的逆作用。真实的数值代码应该解线性系统

\[
S_t Z = C P_t^-
\]

而不是显式求逆——求逆比求解慢且数值不稳定。

协方差更新的简式 \(P_t = P_t^- - K_t C P_t^-\) 在有限精度运算下可能丢失对称性或半正定性。更稳健的 Joseph form 同时从左右两侧应用修正：

\[
P_t = (I - K_t C) P_t^- (I - K_t C)^{\trans} + K_t R K_t^{\trans}.
\]

对于病态问题，平方根 filter 直接传播 Cholesky 因子而非协方差矩阵本身，进一步避免舍入误差导致的负``方差''。数值稳定性不是模型之外的附属问题：负方差或非对称协方差会污染下一次更新中的概率计算，让 filter 在没有任何模型错误的情况下自我毁坏。

还有一种信息形式的 Kalman filter，传播的是协方差矩阵的逆（信息矩阵）和 \(\Omega_t = P_t^{-1} m_t\)。信息 filter 在处理初始不确定性极大（\(P_0 \to \infty\)，即 \(\Omega_0 \to 0\)）的场景时有天然优势——用零信息矩阵比用巨大的协方差矩阵更安全。在传感器网络和多机器人协作估计中，信息形式的加法性质（信息矩阵和信息向量的更新是可加的）使分布式融合变得便捷。

\subsection{非线性与非 Gaussian：闭式递推在此终结}

Kalman filter 的优雅来自一个很窄的前提：线性变换把 Gaussian 映成 Gaussian，联合 Gaussian 的条件分布仍是 Gaussian。这两个性质缺一个，闭式递推就会断裂。

如果动力学或观测是非线性的，Gaussian 经过非线性变换后通常不再是 Gaussian。Extended Kalman filter 的做法是在当前均值处取一阶 Taylor 展开——用局部的线性 \(A_t = \partial f / \partial x\) 和 \(C_t = \partial h / \partial x\) 替代固定的 \(A\) 和 \(C\)。近似的质量取决于：在当前不确定性椭球范围内，真正的非线性函数偏离其线性近似有多远。不确定性很大时，一阶近似可能完全不可靠。一个经典的 EKF 翻车场景是：机器人初始位姿不确定性很大，而观测模型（比如距离和方位的极坐标转换）在远离线性化点处高度非线性——此时 EKF 的线性化可能指向完全错误的更新方向。

Unscented Kalman filter 不取导数，而是用一组确定性选择的 sigma points 穿过非线性函数，再用变换后的样本点重建均值和协方差——它相当于在矩的层面做了至少二阶精度的近似。Particle filter 则完全放弃 Gaussian 假设，用带权的离散粒子表示任意形状的后验分布。

这三种方法之间的层级很清楚：Kalman 假设线性 + Gaussian → 闭式解；EKF/UKF 放开线性 → 近似矩匹配；Particle filter 放开 Gaussian → 任意分布的 Monte Carlo 表示。选择哪一种，取决于你的模型在这条层层放宽的阶梯上站在哪一级——而不是取决于哪一个的代码看起来更现代。

多峰后验是一个特别值得警惕的信号。如果状态确实有两种以上截然不同的可能取值（比如机器人可能在走廊的这一侧或那一侧，但绝不可能在墙里面），用单峰 Gaussian 去近似会把概率质量平均到不可能的区域。这时候即使线性化再精确，被表示族的结构限制也会让 filter 给出一个物理上不可能的估计。问题是表示能力的不足，不是参数精度的不足。Particle filter 也不能解决所有问题——粒子在高维状态空间中的退化（绝大部分粒子的权重趋于零）是它自己的核心困难，需要靠重采样和好 proposal 分布来缓解。

Kalman filter 的力量来源于一组非常具体且严苛的闭包条件：线性映射、Gaussian 噪声与 Markov 状态。看清楚了这些条件，也就看清楚了它在什么情况下是可靠的工具，在什么情况下必须换模型而不只是调协方差参数。
